\chapter{Integral Garis dan Integral Lipat}\label{ch:b2:multint}

Bab ini memperluas pengintegralan dari interval ke kurva dan ke
daerah pada bidang dan pada ruang. \omterm{def:b2:multint:lineint}{Integral garis} mengintegralkan sebuah
\emph{\omterm{def:b2:multint:lineint}{bentuk diferensial}} $P\,\dd x + Q\,\dd y$ sepanjang busur
yang berorientasi; sedangkan integral lipat dua dan lipat tiga mengintegralkan fungsi atas
daerah berdimensi dua dan tiga. Kedua teorinya bertemu pada
\emph{teorema Green--Riemann}, yakni teorema dasar kalkulus
dua dimensi, dan perkakas hitung utamanya di sepanjang bab ini adalah
\emph{rumus penggantian peubah}, yang faktor pemuaiannya adalah
nilai mutlak \omterm{def:b2:linalg:det}{determinan} Jacobi.

\section{Integral garis}

\begin{definition}[Bentuk diferensial; integral garis]\label{def:b2:multint:lineint}
Misalkan $U \subseteq \R^2$ \omterm{def:b2:metric:topology}{terbuka}. Sebuah \emph{bentuk
diferensial}\index{bentuk diferensial} berderajat
$1$ dan berkelas $\mathcal{C}^0$ pada $U$ adalah ungkapan
$\omega = P\,\dd x + Q\,\dd y$ dengan $P, Q \colon U \to \R$
yang \omterm{def:b2:metric:continuity}{kontinu} --- secara formal, sebuah pemetaan \omterm{def:b2:metric:continuity}{kontinu} dari $U$ ke dual
$\R^2$, yakni $\omega(M) = P(M)\,e_1^* + Q(M)\,e_2^*$. Untuk sebuah
busur $\mathcal{C}^1$ $\gamma \colon [a, b] \to U$ dengan $\gamma(t) =
(x(t), y(t))$, \emph{integral garis}\index{integral garis} bagi
$\omega$ sepanjang $\gamma$
adalah
\[
\int_\gamma \omega
= \int_a^b \Bigl(P(\gamma(t))\,x'(t)
+ Q(\gamma(t))\,y'(t)\Bigr)\,\dd t .
\]
Definisinya meluas kata demi kata ke $\R^3$ (dengan bentuk $P\,\dd x +
Q\,\dd y + R\,\dd z$) dan ke busur $\mathcal{C}^1$ sepotong-sepotong (yakni jumlah
atas kepingnya).
\end{definition}

\begin{proposition}[Keinvarianan dan orientasi]\label{prop:b2:multint:lineinv}
Integral garisnya tak berubah di bawah \omterm{def:b2:curves:reparam}{penggantian parameter}
$\mathcal{C}^1$ yang naik, dan berganti tanda di bawah yang turun. Jadi
ia hanya bergantung pada \omterm{def:b2:curves:reparam}{busur geometris} yang \emph{berorientasi}.
\end{proposition}

\begin{proof}
Bila $\theta \colon [c, d] \to [a, b]$ merupakan \omterm{def:b2:curves:reparam}{penggantian parameter} dan
$\tilde\gamma = \gamma \circ \theta$, maka menurut aturan rantai dan
penggantian peubah satu peubah $t = \theta(u)$,
\[
\int_{\tilde\gamma}\omega
= \int_c^d \bigl(P(\gamma(\theta(u)))\,x'(\theta(u))
+ Q(\gamma(\theta(u)))\,y'(\theta(u))\bigr)\,\theta'(u)\,\dd u
= \pm\int_a^b \bigl(Px' + Qy'\bigr)(t)\,\dd t ,
\]
dengan tanda $+$ bila $\theta$ naik (yakni $\theta(c) = a$) dan $-$
bila turun (sebab batasnya bertukar).
\end{proof}

\begin{example}[Usaha sebuah gaya; sirkulasi]\label{ex:b2:multint:work}
Bila $F = (P, Q)$ merupakan medan gaya, maka $\int_\gamma P\dd x + Q\dd y =
\int_a^b \langle F(\gamma(t)), \gamma'(t)\rangle\,\dd t$ adalah
\emph{usaha} $F$ sepanjang $\gamma$. Untuk $\omega = -y\,\dd x +
x\,\dd y$ sepanjang lingkaran satuan berlawanan arah jarum jam $\gamma(t) =
(\cos t, \sin t)$:
\[
\int_\gamma \omega
= \int_0^{2\pi}\bigl((-\sin t)(-\sin t)
+ \cos t\cos t\bigr)\,\dd t = 2\pi ,
\]
yakni dua kali \omterm{def:b2:multint:domain}{luas} yang dilingkupinya --- petunjuk pertama tentang Green--Riemann.
\end{example}

\begin{example}[Satu integral, dua parameterisasi, satu
jebakan tanda]\label{ex:b2:multint:semicircle}
Hitunglah $\int_\gamma x\,\dd y$ sepanjang setengah lingkaran satuan
atasnya dari $(1, 0)$ ke $(-1, 0)$. Dengan $\gamma(t) =
(\cos t, \sin t)$ dan $t \in \intcc0\pi$:
\[
\int_0^\pi\cos t\cdot\cos t\,\dd t = \frac\pi2 .
\]
Sedangkan dengan parameterisasi grafiknya $x \mapsto (x, \sqrt{1 -
x^2})$, dengan $x$ dari $1$ ke $-1$ (perhatikan arahnya!):
\[
\int_1^{-1}x\cdot\frac{-x}{\sqrt{1 - x^2}}\,\dd x
= \int_{-1}^{1}\frac{x^2}{\sqrt{1 - x^2}}\,\dd x
= \frac\pi2
\]
(sebab $x = \sin u$ menyusutkannya menjadi \omterm{pb:b2:multint:1}{integral Wallis}). Jadi nilainya sama,
sebagaimana dijamin \cref{prop:b2:multint:lineinv} --- tetapi hanya
karena kedua jalannya berangkat dari $(1,0)$ ke $(-1,0)$; sebab membalik
perjalanannya membalik tandanya. Adapun menutup lintasannya sepanjang sumbu-$x$
(tempat $\dd y = 0$) tak menambahkan apa pun, dan totalnya $\frac\pi2$
adalah \omterm{def:b2:multint:domain}{luas} setengah cakramnya: yakni wujud pertama bagi
rumus \omterm{def:b2:multint:domain}{luas} perbatasan Green--Riemann di bawah.
\end{example}

\begin{definition}[Bentuk eksak dan bentuk tertutup]\label{def:b2:multint:exact}
Bentuk $\omega = P\,\dd x + Q\,\dd y$ berkelas $\mathcal{C}^0$
disebut \emph{eksak}\index{bentuk eksak} pada $U$ bila ada $f \in \mathcal{C}^1(U)$
(yakni sebuah \emph{potensial}\index{potensial}) dengan $\omega = \dd f$, yakni $P = f_x$ dan $Q
= f_y$. Sedangkan bentuk $\mathcal{C}^1$ disebut \emph{tertutup}\index{bentuk tertutup} bila $P_y = Q_x$ pada
$U$.
\end{definition}

\begin{theorem}[Teorema dasar bagi integral
garis]\label{thm:b2:multint:ftc}
Bila $\omega = \dd f$ bersifat \omterm{def:b2:multint:exact}{eksak} dan $\gamma$ merupakan busur
$\mathcal{C}^1$ sepotong-sepotong di $U$ dari $A$ ke $B$, maka
\[
\int_\gamma \omega = f(B) - f(A) .
\]
Khususnya integral \omterm{def:b2:multint:exact}{bentuk eksak} sepanjang sembarang busur yang tertutup
bernilai nol, dan setiap bentuk $\mathcal{C}^1$ yang \omterm{def:b2:multint:exact}{eksak} bersifat tertutup.
\end{theorem}

\begin{proof}
Berlaku $\frac{\dd}{\dd t}f(\gamma(t)) = f_x(\gamma(t))x'(t) +
f_y(\gamma(t))y'(t)$ menurut aturan rantai
(\cref{ch:b2:diffcalc}), jadi integran pada
\cref{def:b2:multint:lineint} merupakan turunan $t \mapsto
f(\gamma(t))$, sehingga teorema dasar kalkulusnya memberi
hasilnya pada tiap kepingnya; dan nilai antaranya berteleskop. Adapun ketertutupan
bentuk $\mathcal{C}^1$ yang \omterm{def:b2:multint:exact}{eksak} tak lain teorema Schwarz: $P_y =
f_{xy} = f_{yx} = Q_x$.
\end{proof}

\begin{example}[Membangun ulang sebuah potensial]\label{ex:b2:multint:potential}
Misalkan $\omega = y\,\eu^{xy}\,\dd x + (x\,\eu^{xy} + 2y)\,\dd y$
pada $\R^2$. Ia tertutup: sebab kedua turunan silangnya sama dengan
$\eu^{xy}(1 + xy)$. Untuk mencari \omterm{def:b2:multint:exact}{potensialnya}, integralkan $P$ dalam
$x$ dengan $y$ yang tetap:
\[
f(x, y) = \int y\,\eu^{xy}\,\dd x = \eu^{xy} + c(y),
\]
lalu sesuaikan $c$ dengan mencocokkan $f_y$: sebab $x\,\eu^{xy} + c'(y) =
x\,\eu^{xy} + 2y$ memberi $c(y) = y^2$. Jadi $f(x,y) = \eu^{xy} +
y^2$, dan untuk sembarang busur $\mathcal C^1$ sepotong-sepotong dari $(0,0)$
ke $(1,1)$,
\[
\int_\gamma\omega = f(1,1) - f(0,0) = (\eu + 1) - 1 = \eu ,
\]
yang tak bergantung pada lintasannya --- jadi resep dua langkahnya (yakni integralkan
dalam $x$, lalu koreksi dalam $y$) merupakan konvers praktis
\cref{thm:b2:multint:ftc} pada daerah tempat \omterm{def:b2:multint:exact}{bentuk tertutupnya}
bersifat \omterm{def:b2:multint:exact}{eksak}.
\end{example}

\begin{example}[Tertutup tak mengakibatkan eksak]\label{ex:b2:multint:angleform}
Pada $U = \R^2 \setminus \{0\}$, \emph{bentuk sudutnya}
\[
\omega = \frac{-y\,\dd x + x\,\dd y}{x^2 + y^2}
\]
bersifat tertutup (lewat perhitungan langsung: sebab $P_y$ maupun $Q_x$ sama dengan
$\frac{y^2 - x^2}{(x^2+y^2)^2}$), tetapi integralnya sepanjang lingkaran
satuannya adalah $2\pi \neq 0$ (yakni perhitungan yang sama seperti pada
\cref{ex:b2:multint:work}, dibagi $1$): jadi $\omega$ tak \omterm{def:b2:multint:exact}{eksak}
pada $U$. Secara lokal, $\omega = \dd\theta$ untuk sebuah penentuan
$\theta$ bagi sudut kutubnya; jadi kegagalannya bersifat global --- sebab sudutnya
tak dapat didefinisikan secara \omterm{def:b2:metric:continuity}{kontinu} mengelilingi tusukannya. Adapun pada daerah
tanpa lubang patologinya lenyap: sebab pada himpunan \omterm{def:b2:metric:topology}{terbuka} yang \emph{berbentuk
bintang}, setiap bentuk $\mathcal{C}^1$ yang tertutup bersifat \omterm{def:b2:multint:exact}{eksak}
(yakni lema Poincaré, \cref{exo:b2:multint:8}).
\end{example}

\section{Integral lipat dua}

Kita menerima begitu saja teori integral Riemann satu
peubah (yakni jilid Tahun ke-1, dan \cref{ch:b2:integration}) lalu mensketsakan
versi dua peubahnya. Sebuah fungsi $f$ yang \omterm{def:b2:metric:continuity}{kontinu} pada persegi \omterm{def:b2:curves:length}{panjang}
$R = [a, b] \times [c, d]$ punya integral lipat dua $\iint_R f$, yang
didefinisikan lewat jumlah Riemann atas kisi persis seperti pada satu peubah, dan
dihitung lewat iterasi:

\begin{theorem}[Fubini pada sebuah persegi panjang]\label{thm:b2:multint:fubini}
Untuk $f$ yang \omterm{def:b2:metric:continuity}{kontinu} pada $R = [a,b] \times [c,d]$,
\[
\iint_R f
= \int_a^b \Bigl(\int_c^d f(x, y)\,\dd y\Bigr)\dd x
= \int_c^d \Bigl(\int_a^b f(x, y)\,\dd x\Bigr)\dd y .
\]
\end{theorem}

\begin{proof}
Tetapkan $F(x) = \int_c^d f(x, y)\,\dd y$. Kekontinuan seragam $f$ pada
$R$ yang \omterm{def:b2:metric:compact}{kompak} membuat $F$ \omterm{def:b2:metric:continuity}{kontinu} (lewat taksiran yang terdominasi:
$\abs{F(x) - F(x')} \leq (d - c)\sup_y\abs{f(x,y) - f(x',y)}$).
Kini bagilah $[a,b]$ dan $[c,d]$ menjadi $n$ bagian yang sama, sehingga memberi
kisi sel $R_{ij}$ berluas $\Delta x\,\Delta y$. Pada tiap selnya,
$\inf_{R_{ij}} f \cdot \Delta x \Delta y \leq
\int_{x_{i-1}}^{x_i}\int_{y_{j-1}}^{y_j} f(x,y)\,\dd y\,\dd x \leq
\sup_{R_{ij}} f \cdot \Delta x \Delta y$ berkat kemonotonan
integral satu peubahnya (yang diterapkan dua kali). Lalu setelah dijumlahkan atas selnya,
integral teriterasinya $\int_a^b F$ terjepit di antara jumlah Riemann bawah dan
atas kisinya; dan berkat kekontinuan \omterm{def:b2:funcseq:def}{seragamnya} kedua jumlahnya
konvergen ke nilai bersama yang mendefinisikan $\iint_R f$ ketika $n \to
\infty$. Hujah yang sama berlaku dengan peran $x$ dan $y$
dipertukarkan, jadi kedua integral teriterasinya sama dengan $\iint_R f$.
\end{proof}

\begin{remark}
Kekontinuan pada persegi \omterm{def:b2:curves:length}{panjang} yang \omterm{def:b2:metric:compact}{kompak} sungguh bekerja pada
bukti Fubini: sebab ia menyediakan kekontinuan seragam yang
menjepit jumlah Riemannnya. Adapun untuk integran yang lebih liar
pernyataannya sungguh gagal --- sebab ada fungsi yang kedua
integral teriterasinya ada dan berbeda. Teorema umum yang
jujur, dengan keterintegralan sebagai satu-satunya hipotesisnya, adalah
teorema Fubini bagi integral Lebesgue, yang dibuktikan pada
jilid Tahun ke-3; sedangkan segalanya pada bab ini tinggal di latar
yang \omterm{def:b2:metric:continuity}{kontinu}, tempat bukti dasar di atas
\omterm{def:b2:metric:complete}{lengkap}.
\end{remark}

\begin{definition}[Daerah elementer]\label{def:b2:multint:domain}
Sebuah daerah $D \subseteq \R^2$ disebut \emph{elementer-$y$} bila
\[
D = \{(x, y) : a \leq x \leq b,\
\varphi_1(x) \leq y \leq \varphi_2(x)\}
\]
dengan $\varphi_1 \leq \varphi_2$ yang \omterm{def:b2:metric:continuity}{kontinu} pada $[a,b]$
(dan elementer-$x$: secara simetris). Untuk $f$ yang \omterm{def:b2:metric:continuity}{kontinu} pada
sebuah daerah elementer-$y$ $D$,
\[
\iint_D f
= \int_a^b\Bigl(
\int_{\varphi_1(x)}^{\varphi_2(x)} f(x,y)\,\dd y\Bigr)\dd x ,
\]
lalu kita periksa (lewat memperluas $f$ dengan hujah hampiran, atau
lewat pembagian) bahwa ketika $D$ elementer pada kedua arahnya, kedua
integral teriterasinya sepakat. Adapun daerah yang terpotong menjadi berhingga banyak
keping elementer ditangani lewat keaditifannya, dan
\emph{luas}\index{luas!sebuah daerah bidang} $D$ adalah $\operatorname
{Area}(D) = \iint_D 1$.
\end{definition}

\begin{example}\label{ex:b2:multint:triangle}
Pada segitiga $D = \{0 \leq x \leq 1,\ 0 \leq y \leq x\}$:
\[
\iint_D xy \,\dd x\,\dd y
= \int_0^1 x\Bigl(\int_0^x y\,\dd y\Bigr)\dd x
= \int_0^1 x\cdot\frac{x^2}{2}\,\dd x = \frac18 .
\]
Menukar urutannya ($x$ dari $y$ ke $1$): $\int_0^1
y\bigl(\int_y^1 x\,\dd x\bigr)\dd y = \int_0^1
y\,\frac{1 - y^2}{2}\,\dd y = \frac18$ --- jadi nilai yang sama lewat
perhitungan yang berbeda: sehingga memilih urutan pengintegralannya dengan baik adalah separuh
keterampilannya.
\end{example}

\begin{example}[Ketika hanya satu urutan yang berhasil]\label{ex:b2:multint:swaporder}
Hitunglah $I = \displaystyle\int_0^1\!\!\int_x^1
\eu^{y^2}\,\dd y\,\dd x$. Sebagaimana tertulis, integral dalamnya
$\int\eu^{y^2}\dd y$ tak punya antiturunan yang dasar: jadi
perhitungannya macet. Tetapi daerahnya adalah segitiga $0 \leq
x \leq y \leq 1$, yang elementer pada kedua arahnya; jadi setelah
urutannya ditukar,
\[
I = \int_0^1\!\!\int_0^y \eu^{y^2}\,\dd x\,\dd y
= \int_0^1 y\,\eu^{y^2}\,\dd y
= \Bigl[\tfrac12\eu^{y^2}\Bigr]_0^1 = \frac{\eu - 1}{2} .
\]
Sebab peubah dalamnya $x$ tak muncul di mana pun pada integrannya, jadi
mengintegralkannya lebih dulu menghasilkan persis faktor $y$ yang
membuat integral luarnya seketika. Pelajarannya: Fubini bukan sekadar
izin untuk beriterasi --- ia izin untuk \emph{memilih},
dan urutan yang tepat dapat mengubah integral yang mustahil menjadi
satu baris. Jadi sketsakanlah selalu daerahnya lalu bacalah kedua
pemeriannya sebelum memulai.
\end{example}

\begin{theorem}[Penggantian peubah]\label{thm:b2:multint:changevar}
Misalkan $\Phi \colon U' \to U$ difeomorfisma $\mathcal{C}^1$
antara himpunan \omterm{def:b2:metric:topology}{terbuka} $\R^2$, misalkan $K \subseteq U$ daerah \omterm{def:b2:metric:compact}{kompak}
yang terpotong menjadi keping elementer dengan $K' = \Phi^{-1}(K)$, dan
misalkan $f$ \omterm{def:b2:metric:continuity}{kontinu} pada $K$. Maka
\[
\iint_K f(x, y)\,\dd x\,\dd y
= \iint_{K'} f\bigl(\Phi(u, v)\bigr)\,
\abs{\det J_\Phi(u, v)}\,\dd u\,\dd v .
\]
\end{theorem}
\admitted

\begin{remark}
Bukti \omterm{def:b2:metric:complete}{lengkapnya} --- yakni menghampiri $\Phi$ lewat \omterm{def:b2:diffcalc:differential}{diferensialnya} pada
kisi yang halus lalu mengendalikan sel perbatasannya --- \omterm{def:b2:curves:length}{panjang} walaupun
tak dalam; dan ia dikerjakan selengkapnya pada teori ukuran Tahun ke-3, sebagai
akibat teori Lebesgue. Adapun heuristiknya adalah gambaran yang
sudah dipakai bagi \omterm{def:b2:multint:domain}{luas} permukaan: sebuah persegi kecil bersisi $\dd u$ di
$(u, v)$ terpetakan, hingga orde pertama, ke jajaran genjang yang direntang
$\Phi_u\,\dd u$ dan $\Phi_v\,\dd v$, yang berluas
$\abs{\det J_\Phi}\,\dd u\,\dd v$
(\cref{lem:b2:surfaces:lagrange}).
\end{remark}

\begin{remark}[Metode: memilih penggantian peubahnya]
Tiga gerak refleks meliputi sebagian besar kasusnya. \emph{Kesimetrian
integrannya}: sebab $x^2 + y^2$ memanggil koordinat kutub, sedangkan struktur hasil kali
memanggil sumbu Kartesiusnya. \emph{Bentuk
perbatasannya}: sebab perbatasan $u(x,y) = c_1$ dan $v(x,y) = c_2$ memohon
koordinat $(u, v)$ itu sendiri, seperti pada contoh daerah
hiperbolik di bawah --- sebab daerahnya lalu menjadi persegi \omterm{def:b2:curves:length}{panjang},
dan itulah seluruh kemenangannya. \emph{Struktur linearnya}:
sebab ungkapan dalam $x + y$ dan $x - y$ mengundang perputaran
$45$ derajat atau sebuah geseran (\cref{ex:b2:multint:affine}). Pada semua
kasusnya, tiga kotak yang harus dicentang sebelum mengintegralkan: bahwa pemetaannya
bijeksi dari daerah barunya ke daerah lamanya; bahwa Jacobinya
dihitung \emph{pada arah yang sungguh dipakai} (jadi baliklah di
akhirnya bila lebih mudah); dan bahwa Jacobinya masuk dengan
nilai mutlaknya.
\end{remark}

\begin{example}[Koordinat kutub]\label{ex:b2:multint:polar}
Pemetaan $\Phi(\rho, \alpha) = (\rho\cos\alpha,\ \rho\sin\alpha)$ punya
\[
J_\Phi = \begin{pmatrix}
\cos\alpha & -\rho\sin\alpha\\
\sin\alpha & \rho\cos\alpha
\end{pmatrix},
\qquad \det J_\Phi = \rho ,
\]
jadi $\dd x\,\dd y = \rho\,\dd\rho\,\dd\alpha$. Untuk cakram $D_R$
berjari-jari $R$:
\[
\iint_{D_R} e^{-(x^2 + y^2)}\,\dd x\,\dd y
= \int_0^{2\pi}\!\!\int_0^R e^{-\rho^2}\rho\,\dd\rho\,\dd\alpha
= \pi\bigl(1 - e^{-R^2}\bigr)
\xrightarrow[R\to\infty]{} \pi .
\]
Lalu membandingkannya dengan persegi $[-R, R]^2$ (yang terjepit di antara
cakram $D_R$ dan $D_{R\sqrt2}$, dengan semua integrannya positif) memberi
$\bigl(\int_{-\infty}^\infty e^{-x^2}\dd x\bigr)^2 = \pi$:
\[
\boxed{\ \int_{-\infty}^{+\infty} e^{-x^2}\,\dd x = \sqrt{\pi}\ }
\]
--- yakni integral Gauss lagi, kini lewat buktinya yang paling terkenal
(bandingkan dengan penurunan satu peubahnya pada
\cref{ch:b2:integration}).
\end{example}

\begin{example}[Penggantian peubah yang afin]\label{ex:b2:multint:affine}
Untuk \omterm{def:b2:affine:subspace}{pemetaan afin} $\Phi(u, v) = M(u, v)^{\mathsf T} + C$ dengan
$M$ yang terbalikkan, Jacobinya adalah matriks tetap $M$:
jadi luasnya dikalikan faktor tetap $\abs{\det M}$
--- sehingga janji yang dibuat pada \cref{ch:b2:affine} kini menjadi teorema.
Ada dua pemakaian seketika. Elips $\frac{x^2}{a^2} +
\frac{y^2}{b^2} \leq 1$ adalah peta cakram satuannya oleh
$(u, v) \mapsto (au, bv)$, jadi luasnya $ab \cdot \pi$ ---
tanpa perhitungan. Lalu untuk integral $f(x + y)$ atas
persegi $K = \intcc01^2$, geseran $\Phi(u, v) = (u - v, v)$
(yang berdeterminan $1$) mengubahnya menjadi integral $f(u)$ atas sebuah
jajaran genjang, yang diiris Fubini pada $u$ yang tetap: jadi dengan $f =
\exp$,
\[
\iint_K \eu^{x+y}\,\dd x\,\dd y
= \Bigl(\int_0^1 \eu^x\,\dd x\Bigr)^2 = (\eu - 1)^2,
\]
sebagaimana ditegaskan struktur hasil kalinya. Jadi memilih koordinat
yang selaras dengan integrannya --- bukan dengan daerahnya --- adalah
separuh keterampilan yang lain.
\end{example}


\begin{example}[Koordinat yang selaras dengan daerah
lengkung]\label{ex:b2:multint:hyperbolicregion}
Misalkan $D$ daerah pada kuadran pertama yang dibatasi
hiperbola $xy = 1$ dan $xy = 3$ serta garis $y = x$ dan
$y = 3x$. Dalam koordinat $u = xy$ dan $v = y/x$ daerahnya
menjadi persegi $\intcc13 \times \intcc13$; lalu setelah dibalik,
\[
x = \sqrt{u/v}, \qquad y = \sqrt{uv},
\qquad
\det J = x_uy_v - x_vy_u = \frac{1}{2v}
\]
(lewat perhitungan dua baris dengan $x = u^{1/2}v^{-1/2}$ dan $y =
u^{1/2}v^{1/2}$). Karena itu
\[
\operatorname{Area}(D)
= \int_1^3\!\!\int_1^3\frac{\dd u\,\dd v}{2v}
= 2\cdot\frac{\ln 3}{2} = \ln 3 \approx 1.10 .
\]
Mencoba mengiris $D$ dalam koordinat Kartesius berarti memotongnya
menjadi tiga keping dengan perbatasan hiperbolik dan linear ---
yang mungkin, tak menyenangkan, dan rawan salah. Adapun pelajarannya mengulang
\cref{ex:b2:multint:affine} dengan kekuatan penuh: bacalah
persamaan perbatasannya, lalu biarkanlah \emph{ia} memilih
koordinatnya; dan Jacobinya lalu mengubah \omterm{def:b2:multint:domain}{luas} sel jaring
lengkungnya, persis seperti $\rho$ pada koordinat
kutubnya.
\end{example}

\begin{example}[Nilai rata-rata]\label{ex:b2:multint:average}
Adapun \emph{nilai rata-rata} $f$ atas sebuah daerah $D$ adalah
$\frac1{\operatorname{Area}(D)}\iint_Df$. Contohnya: jarak rata-rata
dari pusatnya bagi sebuah titik yang dipilih seragam pada
cakram berjari-jari $R$ adalah
\[
\frac{1}{\pi R^2}\int_0^{2\pi}\!\!\int_0^R
\rho\cdot\rho\,\dd\rho\,\dd\alpha
= \frac{2\pi R^3/3}{\pi R^2} = \frac{2R}3 ,
\]
bukan $R/2$: sebab luas yang seragam menaruh lebih banyak massa pada jari-jari besar (yakni
anulus pada jari-jari $\rho$ berbobot sebanding dengan $\rho$),
jadi rata-ratanya duduk melampaui titik tengahnya. Adapun mendapatkan
faktor ini secara benar persis merupakan Jacobi kutubnya yang bekerja, dan
pembobotan yang sama menjelaskan mengapa titik berat $\bar z = 3R/8$ pada
setengah bola yang dihitung kelak di bab ini bukan $R/2$.
\end{example}

\section{Teorema Green--Riemann}

\begin{theorem}[Green--Riemann]\label{thm:b2:multint:green}
Misalkan $K \subseteq \R^2$ daerah \omterm{def:b2:metric:compact}{kompak} yang elementer pada
kedua arahnya (atau gabungan berhingga daerah demikian yang dilem sepanjang ruas),
dengan perbatasan $\partial K$ berupa kurva tertutup $\mathcal{C}^1$
sepotong-sepotong yang diorientasikan \emph{berlawanan arah jarum jam} (jadi daerahnya tinggal di
kirinya). Untuk $P, Q$ berkelas $\mathcal{C}^1$ pada sebuah persekitaran
$K$:
\[
\oint_{\partial K} P\,\dd x + Q\,\dd y
= \iint_K \Bigl(\frac{\partial Q}{\partial x}
- \frac{\partial P}{\partial y}\Bigr)\,\dd x\,\dd y .
\]
\end{theorem}

\begin{proof}
Pertama, kedua ruasnya bersifat aditif ketika $K$ dipotong sepanjang sebuah
ruas menjadi dua keping $K_1, K_2$: sebab integral lipat duanya bertambah
berkat keaditifan $\iint$; sedangkan untuk integral perbatasannya,
perbatasan berlawanan arah jarum jam milik $K_1$ dan $K_2$ masing-masing menempuh
potongan dalamnya sekali, pada arah yang \emph{berlawanan}, jadi pada
jumlahnya
\[
\oint_{\partial K_1} + \oint_{\partial K_2}
= \oint_{\partial K} + (\text{potongannya, kedua arah})
= \oint_{\partial K},
\]
kedua lintasan sepanjang potongannya saling meniadakan
(\cref{prop:b2:multint:lineinv}) sehingga hanya perbatasan luarnya
bertahan. Lalu setelah berhingga banyak potongan diiterasikan, cukuplah menangani
daerah yang elementer. Kita membuktikan $\oint P\,\dd x = -\iint_K P_y$ pada
daerah elementer-$y$ $D = \{a \leq x \leq b,\ \varphi_1(x) \leq y
\leq \varphi_2(x)\}$; sedangkan kesamaan $\oint Q\,\dd y = \iint_K Q_x$
bersifat simetris (yakni elementer-$x$), dan teoremanya adalah jumlah keduanya.

Hitunglah integral lipat duanya lewat Fubini dan teorema
dasar satu peubahnya:
\[
\iint_D \frac{\partial P}{\partial y}\,\dd x\,\dd y
= \int_a^b \bigl(P(x, \varphi_2(x)) - P(x, \varphi_1(x))\bigr)
\,\dd x .
\]
Kini perbatasan $D$ yang berlawanan arah jarum jam terdiri atas: grafik
bawahnya $y = \varphi_1(x)$ yang ditempuh dari kiri ke kanan, ruas
tegak kanannya $x = b$ (ke atas), grafik atasnya $y =
\varphi_2(x)$ yang ditempuh \emph{dari kanan ke kiri}, dan ruas tegak
kirinya $x = a$ (ke bawah). Sepanjang ruas tegaknya $x$ bersifat
tetap, jadi keduanya menyumbang $0$ pada $\oint P\,\dd x$; sedangkan grafiknya,
yang diparameterkan oleh $x$, memberi
\[
\oint_{\partial D} P\,\dd x
= \int_a^b P(x, \varphi_1(x))\,\dd x
- \int_a^b P(x, \varphi_2(x))\,\dd x
= -\iint_D \frac{\partial P}{\partial y}\,\dd x\,\dd y .
\qedhere
\]
\end{proof}

\begin{corollary}[Luas lewat perbatasannya]\label{cor:b2:multint:area}
Di bawah hipotesis \cref{thm:b2:multint:green},
\[
\operatorname{Area}(K)
= \oint_{\partial K} x\,\dd y
= -\oint_{\partial K} y\,\dd x
= \frac12\oint_{\partial K} x\,\dd y - y\,\dd x .
\]
\end{corollary}

\begin{proof}
Terapkan Green--Riemann pada $(P, Q) = (0, x)$, $(-y, 0)$ dan
$\frac12(-y, x)$: sebab tiap kalinya $Q_x - P_y = 1$.
\end{proof}

\begin{remark}[Memilih di antara ketiga rumus luasnya]
Ketiga rumus perbatasannya setara, tetapi tak dapat saling dipertukarkan dalam
praktik. Pakailah $\oint x\,\dd y$ ketika parameterisasinya membuat
$\dd y$ sederhana (yakni grafik atas sumbu-$y$), lalu $-\oint y\,\dd x$
secara simetris, dan setengah jumlah yang simetris ketika
parameterisasinya memperlakukan $x$ dan $y$ secara merata --- sebab bagi elipsnya
ia menghasilkan integran yang tetap, tanpa pelinearan
trigonometri sama sekali. Adapun pada perbatasan poligonal setengah jumlahnya
menjadi rumus tali sepatu pada \cref{exo:b2:multint:12},
yakni algoritma para juru ukur. Dan ketika perbatasannya ditempuh
searah jarum jam oleh parameterisasi yang diberikan, ketiga rumusnya
mengembalikan \emph{minus} luasnya: jadi hasil yang negatif bukanlah
kesalahan hitung melainkan laporan orientasi --- jadi baliklah
tandanya, atau parameterisasinya.
\end{remark}

\begin{example}[Luas elipsnya]\label{ex:b2:multint:ellipsearea}
Untuk $x = a\cos t$, $y = b\sin t$, dengan $t \in [0, 2\pi]$:
\[
\operatorname{Area}
= \frac12\int_0^{2\pi}\bigl(a\cos t \cdot b\cos t
- b\sin t\cdot(-a\sin t)\bigr)\,\dd t
= \frac{ab}{2}\int_0^{2\pi}\dd t = \pi ab .
\]
\end{example}

\begin{example}[Green--Riemann sebagai
pemeriksaan silang]\label{ex:b2:multint:greencheck}
Ambillah $P = -y^3$ dan $Q = x^3$ pada cakram satuan tertutup $D$.
Sisi perbatasannya, dengan $\gamma(t) = (\cos t, \sin t)$:
\[
\oint_{\partial D}P\,\dd x + Q\,\dd y
= \int_0^{2\pi}\bigl(\sin^4 t + \cos^4 t\bigr)\dd t
= 2\pi\cdot\Bigl(\frac38 + \frac38\Bigr) = \frac{3\pi}2 ,
\]
lewat pelinearan (sebab $\sin^4 + \cos^4 = \tfrac34 +
\tfrac14\cos4t$). Adapun sisi dalamnya:
\[
\iint_D(Q_x - P_y)\,\dd x\,\dd y
= \iint_D 3(x^2 + y^2)\,\dd x\,\dd y
= 3\int_0^{2\pi}\!\!\int_0^1\rho^3\,\dd\rho\,\dd\alpha
= \frac{3\pi}2 .
\]
Jadi bilangan yang sama lewat dua perhitungan yang sangat berbeda --- dan itulah
pemakaian praktisnya: sisi mana pun pada kesamaan Green yang
lebih mudah menjadi perhitungannya, sedangkan yang lain menjadi pemeriksaannya. Untuk
sirkulasi medan polinomial mengelilingi kurva tertutup,
integral lipat duanya hampir selalu sisi yang mudah.
\end{example}

\begin{remark}
Green--Riemann menjelaskan \cref{ex:b2:multint:angleform}: sebab untuk
\omterm{def:b2:multint:exact}{bentuk tertutup} (dengan $Q_x = P_y$), integral mengelilingi perbatasan sembarang
daerah yang termuat di $U$ lenyap. Adapun bentuk sudutnya gagal menjadi \omterm{def:b2:multint:exact}{eksak}
hanya karena tusukan di titik asalnya mencegah cakram yang dibatasi
lingkaran satuannya terletak di dalam $U$ --- jadi \omterm{def:b2:multint:lineint}{integral garis}
\omterm{def:b2:multint:exact}{bentuk tertutup} mendeteksi lubang daerahnya. (Bila didorong lebih jauh,
pengamatan ini menjadi kohomologi de Rham.)
\end{remark}

\section{Integral lipat tiga}

Teorinya meluas ke tiga peubah tanpa gagasan baru: sebab Fubini
menyusutkan $\iiint$ menjadi tiga integral satu peubah, yaitu entah lewat
\emph{pengirisan}: $\iiint_K f = \int\bigl(\iint_{K_z}
f\bigr)\dd z$ atas irisan mendatarnya $K_z$, atau lewat
\emph{penumpukan}: yakni mengintegralkan dalam $z$ lebih dulu sepanjang tongkat tegaknya),
dan rumus penggantian peubahnya berlaku dengan Jacobi
$3 \times 3$.

\begin{example}[Koordinat silinder dan koordinat bola]\label{ex:b2:multint:spherical}
\emph{Silinder} $(x, y, z) = (\rho\cos\alpha, \rho\sin\alpha,
z)$: berlaku $\dd x\,\dd y\,\dd z = \rho\,\dd\rho\,\dd\alpha\,\dd z$.
\emph{Bola} $(x, y, z) = (r\cos\theta\cos\varphi,\
r\sin\theta\cos\varphi,\ r\sin\varphi)$ (dengan $\theta$ bujurnya dan
$\varphi \in [-\frac\pi2, \frac\pi2]$ lintangnya): lalu menguraikan \omterm{def:b2:linalg:det}{determinan}
$3 \times 3$-nya sepanjang baris terakhirnya,
\[
\det J = r^2\cos\varphi ,
\qquad
\dd x\,\dd y\,\dd z
= r^2\cos\varphi\;\dd r\,\dd\theta\,\dd\varphi .
\]
Adapun \omterm{pb:b2:multint:1}{volume bola} berjari-jari $R$:
\[
V = \int_0^R\!\!\int_0^{2\pi}\!\!\int_{-\pi/2}^{\pi/2}
r^2\cos\varphi\;\dd\varphi\,\dd\theta\,\dd r
= \frac{R^3}{3}\cdot 2\pi \cdot 2
= \boxed{\frac43\pi R^3} ,
\]
sehingga akhirnya melunasi rumus yang diterima tanpa bukti pada bab volume
buku terdahulu.
\end{example}

\begin{example}[Tetrahedronnya, dua kali]
\label{ex:b2:multint:tetra}
Volume $T = \{x, y, z \geq 0,\ x + y + z \leq 1\}$, lewat
penumpukan: untuk $(x, y)$ yang tetap pada segitiga $x + y \leq 1$,
$z$ berjalan atas $\intcc0{1 - x - y}$, jadi
\[
V = \int_0^1\!\!\int_0^{1-x}(1 - x - y)\,\dd y\,\dd x
= \int_0^1\frac{(1 - x)^2}{2}\,\dd x = \frac16 .
\]
Sedangkan lewat pengirisan: irisannya pada ketinggian $z$ adalah segitiga $\{x, y
\geq 0,\ x + y \leq 1 - z\}$ yang berluas $\frac{(1-z)^2}2$, jadi
$V = \int_0^1\frac{(1-z)^2}2\,\dd z = \frac16$ lagi --- sebab
kedua perhitungannya adalah integral yang sama dalam urutan yang
berbeda, dan itulah seluruh klaim Fubini. Adapun nilai $\frac16 =
\frac13\cdot\frac12\cdot1$ merupakan rumus kerucutnya
(\cref{ex:b2:multint:cone}) dengan alas segitiga, dan
versi berdimensi $n$-nya $1/n!$ dibuktikan persis lewat pengirisan
ini pada soal akhir pekan.
\end{example}

\begin{example}[Titik berat setengah bola]\label{ex:b2:multint:centroid}
Untuk setengah bola atas $H$ berjari-jari $R$ (dengan $z \geq 0$),
ketinggian titik beratnya adalah $\bar z = \frac1{V}\iiint_H z$, dengan $V =
\frac23\pi R^3$. Dalam koordinat bola (dengan $z =
r\sin\varphi$ dan $\varphi \in \intcc0{\pi/2}$):
\[
\iiint_H z
= \int_0^R r^3\,\dd r\int_0^{2\pi}\dd\theta
\int_0^{\pi/2}\sin\varphi\cos\varphi\,\dd\varphi
= \frac{R^4}4\cdot2\pi\cdot\frac12 = \frac{\pi R^4}4 ,
\]
jadi
\[
\bar z = \frac{\pi R^4/4}{2\pi R^3/3} = \frac{3R}8 :
\]
titik seimbang sebuah setengah bola pejal duduk pada tiga per delapan
jari-jarinya di atas muka datarnya --- yakni di bawah setengah tingginya
$R/2$, sebagaimana mestinya, karena pejalnya lebih gemuk di dekat alasnya.
Adapun setiap perhitungan titik berat berbentuk demikian: satu integral
momen, satu volume, satu nisbah, dan satu pemeriksaan kemasukakalan
terhadap geometrinya.
\end{example}

\begin{example}[Volume lewat pengirisan: kerucutnya]\label{ex:b2:multint:cone}
Sebuah kerucut beralas \omterm{def:b2:multint:domain}{luas} $A$ dan bertinggi $h$ (dengan puncaknya di atas dan alasnya di $z = 0$):
irisannya pada ketinggian $z$ adalah alasnya yang diskalakan dengan faktor $(1 -
z/h)$, sehingga berluas $A(1 - z/h)^2$. Karena itu
\[
V = \int_0^h A\Bigl(1 - \frac zh\Bigr)^2\dd z = \frac{Ah}{3} :
\]
yakni sepertiga pada rumus sekolahnya, yang berlaku bagi \emph{sembarang}
bentuk alas --- sebab pengirisannya mengubahnya menjadi integral sebuah kuadrat.
\end{example}

\begin{example}[Ambang keterintegralan pada
bidang]\label{ex:b2:multint:threshold}
Untuk $\alpha > 0$ yang mana $\iint_{D}\rho^{-\alpha}\,\dd
x\,\dd y$ konvergen pada cakram satuan yang tertusuk $D$ (yakni sebagai limit
atas anulus $\varepsilon \leq \rho \leq 1$)? Dalam koordinat
kutub,
\[
\int_0^{2\pi}\!\!\int_\varepsilon^1\rho^{-\alpha}\,
\rho\,\dd\rho\,\dd\alpha
= 2\pi\int_\varepsilon^1\rho^{1-\alpha}\,\dd\rho ,
\]
yang konvergen ketika $\varepsilon \to 0$ bila dan hanya bila $1 - \alpha > -1$,
yakni $\alpha < 2$: jadi pada dimensi $2$ eksponen kesingularan
kritisnya adalah dimensinya sendiri, sebab $\rho$ tambahan dari
Jacobinya melunakkan kesingularannya sebanyak satu pangkat. (Demikian pula
$\alpha < 3$ bagi kesingularan titik di ruang, lewat $r^2$.)
Adapun tata buku radial semacam ini adalah cara keterintegralannya
diputuskan sekilas pada kerangka Lebesgue Tahun ke-3 ---
dan itulah sebabnya $\iiint 1/r$ konvergen tanpa susah payah pada
\cref{exo:b2:multint:7}.
\end{example}

\begin{remark}[Jebakan yang sering muncul]
(i) \emph{Orientasi}: sebab \omterm{def:b2:multint:lineint}{integral garis} berganti tanda bersama
arah perjalanannya, dan Green--Riemann menuntut perbatasannya
berlawanan arah jarum jam (dengan daerahnya di kiri); sedangkan untuk daerah berlubang,
perbatasan dalamnya ditempuh \emph{searah jarum jam}. (ii)
\emph{Jacobinya masuk dengan nilai mutlak}: sebab penggantian
peubah tak pernah menghasilkan \omterm{def:b2:multint:domain}{luas} yang negatif, dan melupakan
$\abs{\det}$ biasanya membalik tandanya persis ketika pemetaannya
membalik orientasinya. (iii) \emph{Faktor kutub $\rho$}:
sebab $\dd x\,\dd y = \rho\,\dd\rho\,\dd\alpha$, bukan
$\dd\rho\,\dd\alpha$ --- yakni kesalahan yang paling sering di seluruh
bab ini; dan analisis dimensinya menangkapnya, karena
$\dd\rho\,\dd\alpha$ berdimensi \omterm{def:b2:curves:length}{panjang}, bukan
\omterm{def:b2:multint:domain}{luas}. (iv) \emph{Integral lipat dua yang tak wajar}: sebab limit atas
cakram yang membesar dan atas persegi yang membesar sepakat di sini karena
integrannya positif (lewat penjepitan); sedangkan untuk integran
yang bergantitanda limitnya boleh bergantung pada penghabisannya, dan tak ada
klaim yang dibuat tanpa kekonvergenan \omterm{def:b2:series:def}{mutlak}. (v)
\emph{Daerah lawan integran}: sebab integran hasil kali pada daerah
yang bukan hasil kali \emph{tidak} memfaktorkan integralnya ---
sehingga pemfaktorannya menuntut keduanya, seperti pada persegi
\cref{ex:b2:multint:polar}.
\end{remark}

\begin{remark}[Pandangan ke depan di dalam jilid ini]
Integral Gauss yang dihitung di sini diam-diam ada di mana-mana pada
bab peluang: sebab konstanta $\sqrt\pi$ di dalam
rumus Stirling (\cref{thm:b2:comparison:stirling}) adalah
integral bab ini, dan lewat Stirling ia menetapkan
asimtotik $1/\sqrt{\pi n}$ bagi peluang kepulangan jalan
acaknya pada \cref{ch:b2:proba}. Adapun \omterm{pb:b2:multint:1}{integral Wallis} pada
soal akhir pekan muncul kembali di sana juga, sambil menggerakkan
taksiran binomial pusat yang sama. Pada arah lain, unsur \omterm{def:b2:multint:domain}{luas}
dan volume bab ini merampungkan geometri
\cref{ch:b2:surfaces}, dan rumus Green menghitung ulang
\omterm{def:b2:multint:domain}{luas} \omterm{pb:b2:curves:1}{selubung} pada \cref{ch:b2:curves} (yakni \omterm{pb:b2:curves:1}{astroidnya}, pada
\cref{exo:b2:multint:5}). Jadi satu bab, tiga layanan:
ukuran bagi geometri, konstanta bagi peluang, dan
disiplin penggantian peubah yang dipakai keduanya.
\end{remark}

\section{Latihan}

\begin{exercise}[$\star$]\label{exo:b2:multint:1}
Hitunglah $\int_\gamma y^2\,\dd x + x\,\dd y$ sepanjang: (a) ruas
dari $(0,0)$ ke $(1,1)$; (b) busur parabola $y = x^2$ dari
$(0,0)$ ke $(1,1)$. Apakah bentuknya \omterm{def:b2:multint:exact}{eksak}?
\end{exercise}

\begin{exercise}[$\star$]\label{exo:b2:multint:2}
Tunjukkan bahwa $\omega = (2xy + y^3)\,\dd x + (x^2 + 3xy^2 + 1)\,\dd y$
bersifat tertutup pada $\R^2$, carilah sebuah \omterm{def:b2:multint:exact}{potensialnya}, lalu hitunglah
$\int_\gamma\omega$ sepanjang sembarang busur dari $(0, 0)$ ke $(1, 2)$.
\end{exercise}

\begin{exercise}[$\star$]\label{exo:b2:multint:3}
Hitunglah $\iint_D (x + y)\,\dd x\,\dd y$ dengan $D$ menyatakan daerah
yang dibatasi $y = x^2$ dan $y = x$ (dengan $0 \leq x \leq 1$), pada kedua
urutan pengintegralannya.
\end{exercise}

\begin{exercise}[$\star\star$]\label{exo:b2:multint:4}
Dengan memakai koordinat kutub, hitunglah $\iint_D \frac{\dd x\,\dd y}{(1 +
x^2 + y^2)^2}$ atas seluruh bidangnya (sebagai limit atas cakram), dan
$\iint_{D'} xy\,\dd x\,\dd y$ atas seperempat cakram $D' = \{x, y
\geq 0,\ x^2 + y^2 \leq 1\}$.
\end{exercise}

\begin{exercise}[$\star\star$]\label{exo:b2:multint:5}
Hitunglah \omterm{def:b2:multint:domain}{luas} yang dilingkupi \omterm{pb:b2:curves:1}{astroid} $x = \cos^3 t$, $y =
\sin^3 t$, dengan $t \in [0, 2\pi]$, memakai
\cref{cor:b2:multint:area}. \emph{(Linearkan $\sin^2 t\cos^2 t$.)}
\end{exercise}

\begin{exercise}[$\star\star$]\label{exo:b2:multint:6}
Hitunglah volume benda yang dibatasi di bawah oleh paraboloid $z
= x^2 + y^2$ dan di atas oleh bidang $z = 1$, lewat kedua metodenya:
yakni penumpukan (integralkan $1 - x^2 - y^2$ atas cakram satuannya lewat
koordinat kutub) dan pengirisan (sebab irisan mendatarnya berupa cakram berjari-jari
$\sqrt z$).
\end{exercise}

\begin{exercise}[$\star\star$]\label{exo:b2:multint:7}
(Tarikan gravitasi sebuah bola --- teorema Newton, kasus
khusus) Tunjukkan bahwa volume cangkang bola $a \leq r \leq
b$ adalah $\frac43\pi(b^3 - a^3)$ lalu hitunglah $\iiint_{B}
\frac{\dd x\,\dd y\,\dd z}{r}$ atas bola $B$ berjari-jari $R$
(dengan $r$ menyatakan jarak ke titik asalnya). \emph{(Lewat koordinat
bola.)}
\end{exercise}

\begin{exercise}[$\star\star\star$]\label{exo:b2:multint:8}
(Lema Poincaré, kasus berbentuk bintang) Misalkan $U$ berbentuk bintang
terhadap $0$ (yakni $M \in U \Rightarrow [0, M] \subseteq U$)
dan $\omega = P\dd x + Q\dd y$ sebuah bentuk $\mathcal{C}^1$ yang tertutup pada
$U$. Definisikanlah
\[
f(x, y) = \int_0^1 \bigl(x\,P(tx, ty) + y\,Q(tx, ty)\bigr)\dd t .
\]
Dengan memakai penurunan di bawah tanda integral
(\cref{ch:b2:integration}) dan $P_y = Q_x$, tunjukkanlah bahwa $f_x = P$
dan $f_y = Q$: jadi setiap \omterm{def:b2:multint:exact}{bentuk tertutup} pada himpunan \omterm{def:b2:metric:topology}{terbuka} berbentuk bintang bersifat
\omterm{def:b2:multint:exact}{eksak}.
\end{exercise}

\begin{exercise}[$\star\star\star$]\label{exo:b2:multint:9}
(Integral Dirichlet lewat pengintegralan ganda) Benarkan lalu eksploitasilah
\[
\int_0^\infty\!\!\int_0^\infty e^{-xy}\sin x\;\dd y\,\dd x
\quad\text{lawan}\quad
\int_0^\infty\!\!\int_0^\infty e^{-xy}\sin x\;\dd x\,\dd y
\]
pada $[0, A] \times [0, \infty)$: tunjukkan $\int_0^A \frac{\sin
x}{x}\dd x = \frac\pi2 - \int_0^\infty
e^{-Ay}\frac{y\sin A + \cos A}{1 + y^2}\dd y$ lalu perolehlah
$\int_0^\infty \frac{\sin x}{x}\,\dd x = \frac\pi2$, sambil membandingkannya
dengan bukti \omterm{thm:b2:integration:continuity}{integral berparameter} pada
\cref{ch:b2:integration}.
\end{exercise}

\begin{exercise}[$\star\star\star$]\label{exo:b2:multint:10}
(Ketaksamaan isoperimetrik lewat Wirtinger) Misalkan $\gamma$ kurva
tertutup sederhana $\mathcal{C}^1$ berpanjang $2\pi$, yang diparameterkan lewat \omterm{def:b2:curves:length}{panjang
busur} pada $[0, 2\pi]$, dan melingkupi \omterm{def:b2:multint:domain}{luas} $A$. Dengan memakai
\cref{cor:b2:multint:area}, Parseval dan ketaksamaan Wirtinger
(yakni latihan \cref{ch:b2:fourier}), buktikanlah $A \leq \pi$, dengan
kesamaan bagi lingkarannya. \emph{(Normalkan $\int_0^{2\pi} x(s)\dd s
= 0$; lalu tulislah $2A = \oint x\,\dd y - y\,\dd x$ kemudian batasilah $2A \leq
\int (x^2 + y'^2)$ secara berhati-hati lewat $2A = \int_0^{2\pi}(xy' - yx')\dd
s$ dan $x^2 + y'^2 \geq 2xy'$.)}
\end{exercise}

\begin{exercise}[$\star\star$]\label{exo:b2:multint:11}
(Momen bolanya) Untuk bola $B$ berjari-jari $R$ di
$\R^3$, hitunglah $\iiint_B z^2\,\dd x\,\dd y\,\dd z$ dalam
koordinat bola, lalu turunkan $\iiint_B (x^2 + y^2 +
z^2)\,\dd x\,\dd y\,\dd z$ lewat kesimetriannya. Periksalah silang
yang terakhir terhadap perhitungan cangkangnya $\int_0^R r^2\cdot4\pi
r^2\,\dd r$.
\end{exercise}

\begin{exercise}[$\star\star$]\label{exo:b2:multint:12}
(Rumus tali sepatu) Misalkan $K$ sebuah poligon bertitik sudut $(x_1,
y_1), \dots, (x_m, y_m)$ dalam urutan berlawanan arah jarum jam (dengan indeks
modulo $m$). Turunkan dari \cref{cor:b2:multint:area} bahwa
\[
\operatorname{Area}(K)
= \frac12\sum_{i=1}^m
\bigl(x_iy_{i+1} - x_{i+1}y_i\bigr) ,
\]
lalu periksalah rumusnya pada segitiga $(0,0)$, $(1,0)$,
$(0,1)$.
\end{exercise}

\section{Soal: volume bola pada dimensi $n$}

\begin{problem}[{Soal akhir pekan --- $V_n =
\pi^{n/2}/\Gamma(\frac n2 + 1)$, dan keanehan dimensi
tinggi}]\label{pb:b2:multint:1}
Cakram berluas $\pi$, bola bervolume $\frac43\pi$ --- lalu
apa? Soal ini menghitung \omterm{pb:b2:multint:1}{volume bola} satuan
$\R^n$ untuk setiap $n$, dua kali (yakni lewat rekursi pengirisan yang digerakkan
\omterm{pb:b2:multint:1}{integral Wallis}, lalu lewat fungsi $\Gamma$ dan
integral Gauss pada \cref{ex:b2:multint:polar}), lalu membaca
geometrinya: bahwa volumenya memuncak pada dimensi lima lalu bergegas
ke nol, dan hampir seluruh isi bola berdimensi tinggi bersembunyi pada
cangkang tipis di dekat perbatasannya.\index{volume bola}\index{integral
Wallis} Untuk fungsi \omterm{def:b2:metric:continuity}{kontinu} pada bola $\R^n$,
integralnya dipahami sebagai integral teriterasi $n$ lipat
(yakni mengiris satu koordinat sekali waktu, seperti pada bab ini untuk $n
\leq 3$); dan kita menulis $B_n(R)$ bagi bola tertutup berjari-jari $R$
yang berpusat di $0$, $v_n(R)$ bagi volumenya, serta $V_n =
v_n(1)$, dengan $V_0 = 1$ menurut konvensinya.

\medskip
\noindent\textbf{Bagian I --- Rekursi pengirisannya.}
\begin{enumerate}
  \item Dengan menyulihkan $x_i = Ru_i$ pada masing-masing $n$ integral
        teriterasinya, tunjukkanlah $v_n(R) = V_nR^n$.
  \item Dengan mengiris $B_n(1)$ sepanjang koordinat terakhirnya, tunjukkan
        \[
        V_n = V_{n-1}\int_{-1}^{1}(1 -
        t^2)^{\frac{n-1}2}\,\dd t .
        \]
  \item Dengan $t = \sin\theta$, kenalilah integralnya sebagai
        \omterm{pb:b2:multint:1}{integral Wallis}: yakni $\int_{-1}^1(1 -
        t^2)^{\frac{n-1}2}\dd t = 2W_n$, dengan $W_n =
        \int_0^{\pi/2}\cos^n\theta\,\dd\theta =
        \int_0^{\pi/2}\sin^n\theta\,\dd\theta$.
  \item Buktikan kedua kesamaan Wallisnya (yakni integralkan secara parsial;
        lalu teleskopkan $nW_nW_{n-1}$):
        \[
        W_n = \frac{n-1}nW_{n-2} \quad (n \geq 2),
        \qquad
        W_nW_{n-1} = \frac{\pi}{2n} \quad (n \geq 1).
        \]
\end{enumerate}

\medskip
\noindent\textbf{Bagian II --- Rekursinya diselesaikan.}
\begin{enumerate}[resume]
  \item Gabungkanlah pertanyaan 2--4 menjadi rekursi dua langkahnya
        \[
        V_n = \frac{2\pi}{n}\,V_{n-2}
        \qquad (n \geq 2).
        \]
  \item Turunkan \omterm{def:b2:multint:exact}{bentuk tertutupnya}, untuk $k \geq 0$:
        \[
        V_{2k} = \frac{\pi^k}{k!},
        \qquad
        V_{2k+1} = \frac{2^{k+1}\pi^k}{1\cdot3\cdot5\cdots
        (2k+1)} .
        \]
  \item Tabelkan $V_1, \dots, V_7$ secara numerik. Lalu dengan memakai
        nisbah $V_n/V_{n-2} = 2\pi/n$ beserta nilai $2W_5$
        dan $2W_6$, buktikanlah bahwa barisan $(V_n)$ naik
        sampai maksimumnya $V_5 = \frac{8\pi^2}{15} \approx
        5.26$ lalu turun sejak itu.
  \item Tunjukkan bahwa $V_n \to 0$ lebih cepat daripada sembarang barisan
        geometri, dan bahwa $\sum_{n\geq1} V_n$ konvergen: jadi semua
        bola satuannya bersama-sama bervolume total yang berhingga.
  \item Buktikan kesamaan pembangkitnya
        \[
        \sum_{k\geq0} V_{2k}\,x^{2k} = \eu^{\pi x^2}
        \qquad (x \in \R),
        \]
        lalu turunkan $\sum_{k\geq0}V_{2k} = \eu^\pi \approx
        23.14$.
\end{enumerate}

\medskip
\noindent\textbf{Bagian III --- Rute kedua: $\Gamma$ dan
integral Gauss.}
\begin{enumerate}[resume]
  \item Tunjukkan lewat Fubini (sebab integrannya berupa hasil kali) bahwa
        \[
        I_n = \int_{\R^n}\eu^{-\norm
        x^2}\dd x
        = \Bigl(\int_{-\infty}^{+\infty}
        \eu^{-t^2}\dd t\Bigr)^{\!n} = \pi^{n/2},
        \]
        yakni integral Gauss berdimensi $n$, yang dipahami sebagai
        limit atas kubus $\intcc{-R}{R}^n$.
  \item Ingat kembali $\Gamma(s) = \int_0^\infty t^{s-1}\eu^{-t}\dd
        t$ (\cref{def:b2:integration:gamma}). Lalu dari
        $\Gamma(s+1) = s\,\Gamma(s)$
        (\cref{thm:b2:integration:gammaprops}) dan
        $\Gamma(\tfrac12) = \sqrt\pi$ (yakni sulihkan $t = u^2$
        lalu panggillah integral Gaussnya), hitunglah
        \[
        \Gamma(k + 1) = k!,
        \qquad
        \Gamma\Bigl(k + \frac32\Bigr) =
        \frac{1\cdot3\cdots(2k+1)}{2^{k+1}}\,\sqrt\pi .
        \]
  \item Buktikan, secara induktif lewat rekursi pertanyaan
        5, satu rumus tunggalnya
        \[
        V_n = \frac{\pi^{n/2}}{\Gamma\bigl(\frac n2 +
        1\bigr)} \qquad (n \geq 1),
        \]
        lalu periksalah bahwa ia menghasilkan kembali kedua \omterm{def:b2:multint:exact}{bentuk tertutup} pertanyaan
        6.
  \item Tunjukkan $\int_0^\infty \eu^{-r^2}r^{n-1}\dd r =
        \tfrac12\Gamma\bigl(\tfrac n2\bigr)$ lalu turunkan
        kesamaan
        \[
        I_n = n\,V_n\int_0^\infty \eu^{-r^2}\,r^{n-1}\,\dd r.
        \]
        Tafsirkanlah: bahwa massa Gauss $\R^n$ dikumpulkan
        sepanjang cangkang bola yang
        ``\omterm{def:b2:multint:domain}{luas} berdimensi $(n-1)$''-nya pada jari-jari $r$ adalah
        $nV_nr^{n-1}$ --- dan kedua ruasnya kini terbukti
        secara bebas, jadi tafsirannya tak berharga apa pun.
  \item Tetapkan $s_{n-1} = nV_n$ (yakni \omterm{def:b2:multint:domain}{luas} bola satuan
        $S^{n-1}$, yang konsisten dengan $v_n(R) = \int_0^R
        s_{n-1}r^{n-1}\dd r$). Tabelkanlah $s_0, \dots, s_3$ lalu
        periksalah $s_1 = 2\pi$, $s_2 = 4\pi$, dan $s_3 = 2\pi^2$.
\end{enumerate}

\medskip
\noindent\textbf{Bagian IV --- Dimensi tinggi itu aneh.}
\begin{enumerate}[resume]
  \item Dari rumus Stirling
        (\cref{thm:b2:comparison:stirling}) yang diterapkan pada $k!$,
        tunjukkanlah untuk $n = 2k$ yang genap:
        \[
        V_n \sim \frac{1}{\sqrt{\pi n}}
        \Bigl(\frac{2\pi\eu}{n}\Bigr)^{n/2}
        \qquad (n \to \infty, \ n \text{ genap}),
        \]
        lalu jelaskanlah mengapa batas peluruhan supergeometri yang sama
        meluas ke $n$ yang ganjil lewat rekursinya.
  \item Bola satuannya duduk di kubus $\intcc{-1}1^n$ yang
        bervolume $2^n$. Hitunglah nisbah pengisiannya $V_n/2^n$ untuk
        $n = 2, 3, 10$, lalu tunjukkan bahwa ia menuju $0$: jadi pada dimensi
        tinggi, pada dasarnya seluruh isi kubusnya terletak di
        pojoknya.
  \item Tunjukkan bahwa pecahan $v_n(1)$ yang terletak dalam
        jarak $\varepsilon$ dari bola perbatasannya adalah $1 -
        (1 - \varepsilon)^n \to 1$; lalu secara numerik, berapa
        pecahan bola berdimensi $100$ yang terletak pada
        cangkang luar setebal $1\%$?
  \item Buktikan asimtotik Wallis $W_n \sim
        \sqrt{\dfrac{\pi}{2n}}$ \emph{(lewat kemonotonan
        $(W_n)$, nisbah $W_n/W_{n-2} \to 1$, dan $W_nW_{n-1}
        = \frac\pi{2n}$)}, beserta batas bawahnya $W_n \geq
        \sqrt{\dfrac{\pi}{2(n+1)}}$ untuk setiap $n$.
  \item (Pemusatan pada sebuah lempeng) Adapun pecahan bola satuan
        yang koordinat pertamanya melampaui $\delta$ adalah
        $\int_\delta^1(1 - x^2)^{\frac{n-1}2}\dd x \,\big/\,
        (2W_n)$. Dengan memakai $1 - u \leq \eu^{-u}$ dan batas
        ekornya $\int_\delta^\infty \eu^{-a x^2}\dd x \leq
        \frac{\eu^{-a\delta^2}}{2a\delta}$, tunjukkanlah bahwa pecahan
        ini paling banyak
        \[
        \frac{\eu^{-(n-1)\delta^2/2}}{(n-1)\,\delta}
        \Big/ \sqrt{\frac{2\pi}{n+1}}
        \]
        lalu simpulkan: bahwa untuk $\delta = s/\sqrt{n-1}$, semuanya kecuali
        pecahan $O(\eu^{-s^2/2}/s)$ bagian bolanya terletak pada
        lempeng $\abs{x_1} \leq s/\sqrt{n-1}$. Jadi sebuah bola berdimensi
        tinggi, secara statistik, berupa panekuk tipis pada setiap
        arah sekaligus.
  \item Rangkailah pertanyaan 16--19 menjadi satu paragraf: yakni di mana
        volume $B_n(1)$ duduk (di dekat bola
        perbatasannya, namun di dalam lempeng $O(1/\sqrt n)$ bagi setiap
        hiperbidang lewat pusatnya), dan mengapa kedua
        pernyataan itu tak saling bertentangan.
\end{enumerate}

\medskip
\noindent\textbf{Bagian V --- Benda lain, dan rangkuman.}
\begin{enumerate}[resume]
  \item (Simpleks) Misalkan $\Delta_n = \{x \in \R^n : x_i \geq 0,\
        \sum x_i \leq 1\}$. Tunjukkan lewat pengirisan dan induksi bahwa
        $\operatorname{vol}(\Delta_n) = \frac1{n!}$.
  \item (Politop silang) Turunkan bahwa $C_n = \{x :
        \sum\abs{x_i} \leq 1\}$ bervolume $\frac{2^n}{n!}$,
        lalu periksalah apitan $C_n \subseteq B_n(1)
        \subseteq \intcc{-1}1^n$ pada aras volumenya:
        $\frac{2^n}{n!} \leq V_n \leq 2^n$.
  \item Hitunglah $V_4$ lewat cara ketiga: irislah $\R^4 = \R^2 \times
        \R^2$, integralkan \omterm{def:b2:multint:domain}{luas} cakram-$(z, w)$-nya atas
        cakram-$(x, y)$-nya dalam koordinat kutub, lalu perolehlah
        $V_4 = \frac{\pi^2}2$.
  \item (Monte Carlo dalam kesulitan) Sebuah titik ditarik seragam
        di kubus $\intcc{-1}1^{20}$. Tunjukkan bahwa
        peluangnya mendarat di bola yang terlukis di dalamnya adalah
        $V_{20}/2^{20} \approx 2.5\cdot10^{-8}$, jadi sekitar
        empat puluh juta penarikan diperlukan sebelum sasaran yang pertamanya
        diharapkan: sehingga menaksir $V_n$ lewat pencuplikan penolakan
        runtuh pada dimensi tinggi (yakni kutukan
        dimensinya).
  \item Rangkuman. Dua penurunan yang bebas bertemu di $V_n =
        \pi^{n/2}/\Gamma(\frac n2 + 1)$: daftarkanlah teorema mana
        pada bab ini yang dipakai masing-masingnya (Fubini, penggantian
        peubah, integral Gauss kutubnya), dan masukan
        satu peubah mana (Wallis, $\Gamma$, Stirling).
        Lalu di mana jilid Tahun ke-3 mengerjakan ulang perhitungan ini
        dengan teori Lebesgue, dan apa yang ditambahkannya?

\end{enumerate}
\end{problem}
